% Autoría principal: Roger; arquitectura revisada con todos.
\subsection{Diseño de los tres nodos}
\pendiente{Actualizar el diagrama de FORM-TL-01 con nodo-A, nodo-B, nodo-C y puertos 5000, 5001, 5002, indicando qué proceso actúa como servidor y cliente. No presentar una arquitectura propuesta como ejecutada.}

\subsection{Diagrama de arquitectura (TikZ)}
% Crear figures/arquitectura_tikz.tex desde el diagrama real FORM-TL-01.
\IfFileExists{figures/arquitectura_tikz.tex}{%
 \begin{figure}[htbp]\centering\begingroup
\shorthandoff{<>}
\begin{tikzpicture}[
  node distance=2cm and 3cm,
  every node/.style={align=center},
  box/.style={draw,rounded corners,text width=4cm,minimum height=1.6cm},
  >=Stealth]
\node[box] (a) {nodo-A:5000\\Servidor + cliente\\Reloj A + registro A};
\node[box,right=of a] (b) {nodo-B:5001\\Servidor + cliente\\Reloj B + registro B};
\node[box,below=of a,xshift=3.5cm] (c) {nodo-C:5002\\Servidor + cliente\\Reloj C + registro C};
\draw[->,thick] (a) -- node[above] {solicitudes} (b);
\draw[->,thick] (b) -- node[right] {solicitudes} (c);
\draw[->,thick] (c) -- node[left] {solicitudes} (a);
\end{tikzpicture}
\endgroup

 \caption{Arquitectura validada de los tres nodos (diagrama del equipo).}\label{fig:arquitectura}\end{figure}
}{\pendiente{Crear docs/figures/arquitectura\_tikz.tex con el diagrama real basado en FORM-TL-01.}}

\subsection{Modelo de mensajes y flujo de comunicación}
\pendiente{Mostrar estructura JSON {\ttfamily\{sender, timestamp, message\}}, su equivalente Protocol Buffers, direcciones de mensajes y correlación de solicitud/respuesta.}

\subsection{Ejecución y decisiones de despliegue}
\pendiente{Explicar si se ejecutan procesos en localhost o contenedores Docker, y cómo se evitan conflictos de puertos cuando se comparan los dos mecanismos.}
